\section{Research questions}

There are four key areas of focus that this research seeks to address. The
first two are questions into the nature of biological process.

\begin{itemize}
\item \textbf{Why did \yp{} evolve a biphasic pathogenesis in mammalian hosts?}
Switching from a pro-phagocytic to an anti-phagocytic strategy seems to be
crucial to the overall pathogenesis of \yp{}, and the preceding sections give
reasons why the sequence is forced without saying what it is worth. There is a
simple argument based on the net proliferability both before and after
neutrophils arrive on the scene of infection, which sorts the bacterium's
options into a two-by-two and locates a threshold neutrophil fraction at which
the better option changes; it is set out in
\mextref{con:sec:quadrant}{Chapter~10}. A second argument, of a quite different
kind, keys the same switch to the amount of killing the host can afford per
delivery. And a third strand supports both without deciding either: if \yp{}
replication is only slightly supercritical, then a near-critical branching
process becomes drastically less likely to die out when its initial count is
increased by even a few colony forming units
(\mextref{m:app:smallpops}{Section~2.A}). If we see the within-macrophage
replication as setting up an initial population, this effect offers further
explanatory potential.

\item \textbf{Is there a general advantage to ``bursting'' style transfer in
successive intracellular infection?} N.~Komarova proposed that there could be a
flooding advantage to bursting pathogenesis --- bacteria or virus exiting their
host cells all at once in large groups rather than one at a time or in dribs
and drabs. \Mextref{p:ch:one-cell-to-population}{Chapter~6} and
\mextref{path:ch:pathogenesis}{Chapter~8} give mathematical support to the theory and,
more usefully, say when it holds and when it does not.
\end{itemize}

\noindent
The second two are areas in which we tried to construct new models to describe
particular dynamics.

\begin{itemize}
\item \textbf{An intracellular replacement for the standard model of viral
replication.} The two constants of the BMVR describe a budding cell exactly and
a bursting cell not at all, since for a lytic pathogen release and death are
the same event rather than independent ones. Based on our analysis of the
birth--death--catastrophe process, those constants can be replaced by
quantities derived from the intracellular dynamics, and the question is what
that costs and what it buys.

\item \textbf{Efforts to describe variable rates of evolution with Markov
chains.} The standard birth--death process has been used for decades to model
the emergence of species over geological time periods. It may however be
criticised for being consistent with the outdated doctrine of progressive
gradualism, which holds that evolutionary change is constant and happens
gradually. In 1972 Eldredge and Gould proposed the counter theory of punctuated
equilibria \cite{eldredge1972punctuated}: evolution taking place rapidly over
proportionally short periods, intervened by longer times of relative
morphological stasis. There is also the controversial question of why large
clades appear generally to have experienced an initial period of substantial
diversification, which some attribute to a statistical consequence of standard
birth--death dynamics --- the push of the past, and the pull of the present.
This work presents models which seek to give mathematical basis to the
punctuated-equilibria paradigm, and asks more generally what survives when a
rate is allowed to move. There is potential significance here, too, in
understanding the emergence of \yp{} itself.
\end{itemize}

\noindent
\Mextref{con:sec:questions}{Chapter~10} returns to all four and answers them as
far as the work allows, which in two cases is not all the way.
